% 第五大组共同报告封面；个人详细贡献统一放在分工章。
\begin{titlepage}
\thispagestyle{empty}
\begin{tikzpicture}[remember picture,overlay]
\fill[Navy] (current page.north west) rectangle ([yshift=-12mm]current page.north east);
\end{tikzpicture}
\centering
\vspace*{17mm}
{\sffamily\small\color{Muted} INTELLIGENT ROBOTICS LAB\par}
\vspace{12mm}
{\Large\ReportClass\par}
\vspace{13mm}
{\Huge\bfseries\color{Navy} 基于视觉与多方式跟随的\par 多机器人协同控制\par}
\vspace{8mm}
{\Large\color{Teal} 第五大组共同实验报告\par}
\vspace{8mm}
\rule{.85\linewidth}{.6pt}\par
\vspace{5mm}
\begin{minipage}{.88\linewidth}\centering
TonyPi抱物穿行与标记路线控制\\[3mm]
AR直接跟随\quad 地图定位与避障跟随\\[3mm]
RGB-D实时检测\quad 多车AR队列跟随
\end{minipage}\par
\vspace{10mm}
\begin{minipage}{.88\linewidth}\small
\begin{tabularx}{\linewidth}{p{24mm}Y}
\textbf{第一小组} & 杨开元、林楗棋、曹淇涵、刘小凡\\[4mm]
\textbf{第二小组} & 陈凯华、\LuoMin{}、刘彦廷、许程奕（按仓库投稿）\\[4mm]
\textbf{第三小组} & 成员信息待补充\\[4mm]
\textbf{方式四小组} & 徐杰鸿、王振宇、林琢严、何爽；组号待确认\\[4mm]
\end{tabularx}
\end{minipage}\par
\vfill
\begin{minipage}{.88\linewidth}\small\color{Muted}
报告范围覆盖整个大组的系统设计、各方式实现、组间接口、组内分工与联合实验。各组原始材料分别署明来源；学号及尚缺人员信息在正文登记，未提供的信息不作推定。
\end{minipage}\par
\vspace{8mm}
{\ReportDate\par}
\end{titlepage}
